\section{Rational enumeration at every fixed excess}
\label{sec:fixed-excess-rationality}

The fixed-excess grammar can be strengthened to an effective rational
enumeration theorem.  The proof counts canonical colored incidence
trees, rather than ambiguous grammar derivations.

\begin{theorem}[Rational unlabelled enumeration at fixed excess]
\label{thm:fixed-excess-rationality}
Fix an integer $e\ge0$, and let $\mathcal C_e$ be the isomorphism
classes of two-trees $G$ of order at least six with
$2|\partial G|-|V(G)|-4=e$.  Then
\[
 F_e(z)=\sum_{[G]\in\mathcal C_e}z^{|V(G)|}
\]
is an effectively computable rational function.  In particular,
for each fixed $e$, its exact unlabelled counts eventually satisfy
a constant-coefficient linear recurrence.
\end{theorem}

The restriction to orders at least six avoids the exceptional small
orders in the sharp bound.  Extending a fixed-excess count to all
orders changes its generating function by only a finite polynomial.
No practical complexity bound is asserted for the general algorithm.

\subsection{Finite bounds and an intrinsic pending-root weight}

Put
\[
 C=2e+6,\qquad D=\max(4,2e+3),\qquad L=2e+2.
\]
The structural defect theorem bounds interior degrees by $C$ and
boundary degrees by $D$.  Remove all degree-two graph vertices to
obtain the canonical core $H$.  If $H$ is noncomplete, it is a
two-tree with at most $L$ ears.  Its internal triangle--edge incidence
tree, obtained by deleting private edge-nodes, has at most $L$
leaves.  Every removed vertex is an ear on a core edge $f$; write its
multiplicity as $q_f$.  Since each core endpoint has degree at least
two in $H$,
\[
 0\le q_f\le C-2.
\]
The complete cores $K_2,K_3$ have the already classified orders
$e+4$ or $e+6$ and bounded ear multiplicities, so contribute a
finite polynomial.  We treat noncomplete cores below.

Cut at an internal core edge $uv$, with at least one core triangle
on either side, and defer all ears on that root edge.  Exclude
$u,v$ from the component weight $W$.  Every other vertex receives
its boundary/interior sign under the \emph{pending-root convention}
that the final multiplicity of $uv$ is at least two.  This is not
the raw boundary weight of the isolated component.  Indeed, adding
the first outside triangle can activate a local condition at a page
vertex.  Under the pending convention that condition is already
active: the local formula
\[
 \beta(u)=R_u-\deg u,\quad
 R_u=\max\{2+t_{uv},\ 2+t_{ua}+t_{ub}:
 uab\text{ a triangle},\ t_{ab}\ge2\}
\]
depends on the root edge, for a nonroot vertex, only through the
threshold $t_{uv}\ge2$.  All incident edges and degrees of that
nonroot vertex are closed.  Thus $W$ is intrinsic and independent
of the chosen completion with an active root edge.

Every such component of a graph in $\mathcal C_e$ satisfies
\begin{equation}
 0\le W\le e+6.
 \label{eq:rational-pending-weight}
\end{equation}
To prove the lower bound, add one vertex on the root edge.  Its
multiplicity is now at least two, preserving exactly the pending
nonroot signs just defined.  Every two-tree of order at least three
has total boundary-minus-interior weight at least three, by the
sharp theorem and the direct small-order check.  This completion
has at least four vertices, so a convention also admitting $K_2$
is irrelevant.  Its two former roots and new degree-two vertex
contribute at most three, proving $W\ge0$.

For the upper bound, let $W'$ be the pending nonroot weight of the
complementary component, and $q$ the root ear multiplicity.  With
root signs $s_u,s_v\in\{-1,1\}$, the final graph gives
\[
 e+4=W+W'+q+s_u+s_v.
\]
Since $W'\ge0$, $q\ge0$, and $s_u+s_v\ge-2$, the upper bound in
\eqref{eq:rational-pending-weight} follows.  The argument also
applies to any nonempty proper group of pages at an internal edge.

\subsection{An exact finite-alphabet incidence encoding}

A two-tree has a proper three-coloring unique up to one global
permutation: name the colors on an initial triangle, after which
each added vertex is forced to use its two parents' missing color.
For a named coloring of $H$, color every edge-node of its full
triangle--edge incidence tree by the missing endpoint color.  The
three edge-nodes at a triangle have different colors.  This colored
tree reconstructs the colored graph uniquely.  Root at a triangle,
name its three vertices by color, and traverse the tree; every new
triangle shares a known edge and introduces one new vertex of the
missing color, while its two other edge colors determine the
incidences.  Conversely this procedure reconstructs a two-tree
from every locally valid such colored incidence tree.

Use the internal incidence tree $T$, obtained by deleting all
edge-nodes lying in just one triangle.  Give its nodes finite labels:
\begin{itemize}
\item an internal edge-node carries its missing color and $q_f$;
\item a triangle-node records the subset of edge colors remaining
internal, and the ear multiplicity for each missing private edge
color.
\end{itemize}
Its neighboring edge colors must be precisely that recorded subset.
All deleted private edges can be restored uniquely.  Thus these
labels reconstruct $H$ and then $G$.  Ears on one edge are represented
by their multiplicity, so their permutations are already removed.

The alphabet is finite.  Triangle degrees are at most three, and
edge-node degrees are at most $L$, since $T$ has at most $L$ leaves.
All leaves are triangle-nodes.  Define a node weight by
\[
 \omega(v)=
 \begin{cases}
 1+\sum q_f,&v\text{ a triangle-node, over its private edges},\\
 q_f,&v\text{ an internal edge-node}.
 \end{cases}
\]
Then
\begin{equation}
 |V(G)|=2+\sum_{v\in T}\omega(v).
 \label{eq:rational-node-weight}
\end{equation}
A core with $t$ triangles has $t+2$ vertices, and every ear has been
counted once.  All weights are nonnegative and triangle weights
are positive.  Color permutations preserve them.  Colored decorated
tree isomorphism is exactly colored graph isomorphism; forgetting
colors adds precisely the one global $S_3$ action, with no separate
gluing gauge at each edge.

\subsection{A deterministic recognizer on the canonical tree}

Root $T$ at a designated leaf triangle.  At an internal edge cut,
order the two graph endpoints by their named colors and retain:
the number of inside root triangles; both current endpoint degrees;
both requirements $\max(A,B+t)$, with $A$ capped at $C+1$ and $B$
at $C-1$; the inside leaf count; actual pending weight $W$ in
$\{0,\ldots,e+6\}$; and the finite deferred-ear and color data.
Continuing degrees are at most $C-1$.  The requirement caps preserve
every final comparison at degree at most $C$, since $t\ge2$.
All coordinates are finite; include a rejecting state.

There are canonical deterministic transitions.  A triangle has zero,
one, or two active child core edges, whose ports are distinguished
by color.  The local boundary formula computes the affine root
requirements and the final sign and degree of its third vertex;
labels specify all ear multiplicities.  Add that sign and the ears
to $W$.  At an edge-node combine \emph{one unordered star of all
child triangle branches at once}: root degrees add with subtraction
for the repeated root neighbor, requirements combine by maxima,
and triangle counts, leaves, and nonroot weights add.  This transition
is symmetric in a bounded multiset of child states, retaining the
full multiplicity of repeated branches.  No binary-join history is
counted.  The affine constants can be defined intrinsically as
maxima of constant and root-linear local contributions, rather
than by a choice among equivalent formulas.

Require each finalized retained core vertex to have degree at least
three and satisfy the appropriate boundary/interior cap.  Reject if
any state bound is exceeded.  At the root leaf triangle, close the
remaining edge, account for its ears and the two private edges,
finalize all remaining vertices, and accept exactly at total weight
$e+4$.  The root core vertices must also finish with degree at least
three.

Every member of $\mathcal C_e$ passes under any chosen leaf root,
by the proved bounds and exact local formula.  Conversely an accepted
decorated tree reconstructs a two-tree whose specified core vertices
all have degree at least three and whose attached ears have degree
two.  Its specified core is therefore the canonical pruned core.
All signs are correct and total weight $e+4$ is exactly the desired
excess.  A noncomplete pruned core with all ears retained above degree
two gives order at least six.  Thus the deterministic recognizer
accepts exactly the required graph property.  Intermediate states
may depend on the designated root, but acceptance does not.

\subsection{Finite skeletons and a disjoint language partition}

Suppress all degree-two \emph{incidence nodes} of $T$.  This is
distinct from removing degree-two graph vertices.  The resulting
skeleton $K$ has at most $L$ leaves and $L-2$ branch nodes, hence
at most $2L-2$ nodes.  Choose one named representative of each of
the finitely many abstract skeletons, a designated leaf, and an
orientation of each edge.  A decorated core with this skeleton is
specified by finite labels on the skeleton nodes and one word on
each oriented edge, listing its suppressed degree-two node labels.
Local type and color compatibility are finite conditions.  Chain
reversal is literal word reversal and color action is letterwise.

Along a chain, rooted towards the designated leaf, the recognizer
applies unary maps on its finite state set.  Adjoin rejection to
make every map total.  A chain word has one exact composite map
in a finite transformation monoid.  For any fixed map, its words
form a regular language, and these languages are disjoint.  Endpoint
compatibility can be incorporated in the finite machine.  Once all
skeleton labels and composite chain maps are fixed, acceptance is
decided by the finite remaining computation at skeleton nodes.

The valid configurations on a named skeleton are consequently a
\emph{disjoint finite union of products of regular chain languages},
with fixed skeleton labels.  This exact-map partition avoids
counting a graph once per derivation or state annotation.

The word weight is the sum of the node weights in
\eqref{eq:rational-node-weight}.  Edge letters can have weight zero,
but types alternate and triangle letters have positive weight;
a word of length $h$ has weight at least $\lfloor h/2\rfloor$.
Thus there are finitely many words of each weight.  On a productive
recognizing automaton no zero-weight directed cycle is possible,
and the weighted transfer-matrix expression is a well-defined
rational formal power series.

\subsection{Fixed words under twisted reversal}

We need a finite-language fact that includes reflections.  Let
$\mathcal L$ be a regular language on a finite nonnegatively
weighted alphabet, with finitely many realizations at each weight,
and let $\rho$ be a weight-preserving permutation of its letters.
The generating functions for
\[
 \{w\in\mathcal L:w=\rho(w)\},\qquad
 \{w\in\mathcal L:w=\rho(\operatorname{rev}w)\}
\]
are rational.  The first language restricts to letters fixed by
$\rho$.  For the second, applying the equation twice restricts all
letters to $\operatorname{Fix}(\rho^2)$, on which $\rho$ is an
involution.  Each fixed word is uniquely
\[
 u\rho(\operatorname{rev}u)quad\text{or}\quad
 ua\rho(\operatorname{rev}u),\qquad\rho(a)=a.
\]

Let a deterministic automaton for $\mathcal L$ have initial state
$q_0$, accepting set $F$, and letter transitions $\delta_b$.
While reading $u$, track
\[
 \bigl(\delta(q_0,u),T_{\rho(\operatorname{rev}u)}\bigr),
\]
where $T_v$ is the full state transformation for word $v$.  This
is a finite state.  Appending a letter $b$ updates it by
\begin{equation}
 (p,T)\longmapsto
 \bigl(\delta_b(p),T\circ\delta_{\rho(b)}\bigr).
 \label{eq:rational-half-update}
\end{equation}
For even length accept if $T(p)\in F$.  For an odd center $a$,
accept if $T(\delta_a(p))\in F$.  Therefore the half-words form
regular languages.  Give a half-letter twice its original weight
and add the central weight where present.  This proves rationality,
counting each fixed word once by its unique half-word.  The
alternating-type weight bound also precludes productive zero-weight
cycles here.  Intersections and inverse images under letter
permutations or reversal preserve regularity, so the same argument
applies with finitely many additional chain constraints.

\subsection{The full isomorphism quotient and Burnside count}

For a fixed named skeleton the finite group is
\[
 \Gamma_K=\operatorname{Aut}(K)\times S_3.
\]
The first factor permutes nodes and chains and reverses their words
when it reverses the selected orientation; the second acts globally
on all color labels.  The actions commute and preserve weights.
Their orbits are exactly the uncolored graph isomorphism classes.
An isomorphism preserves the canonical core, its incidence tree,
and its suppression.  On a chain it is determined by the endpoints
and hence preserves or reverses the word; there are no translations
or extra local gauges.  The colored incidence encoding is bijective,
and forgetting the unique named coloring adds precisely $S_3$.
Conversely every such group identification reconstructs an
isomorphism of the graphs.  Ear permutations are already represented
by multiplicity counts.

For $g\in\Gamma_K$, choose an assignment of skeleton labels fixed by $g$,
and one accepting product of the disjoint classes defined by chain maps. On
an orbit of $m$ chains, all words are determined by one representative
$w$ through specified color permutations and reversals.  Pull the
regular conditions on the other chains back to $w$ and intersect
them.  Going around the orbit imposes exactly
\[
 w=\rho(w)\quad\text{or}\quad
 w=\rho(\operatorname{rev}w)
\]
for an accumulated color permutation $\rho$.  The preceding
half-word construction gives a rational representative series.
All $m$ copies have the same weight, so use the substitution
$z\mapsto z^m$.  Different chain orbits are independent once
their finite maps and labels are fixed, so their series multiply.
Summing over the finite disjoint products gives a rational fixed
configuration series for $g$.

The designated recognizer root need not be fixed by $g$: acceptance
is the intrinsic graph property.  Nor must the exact-map partition
be group-invariant; intersecting each disjoint product with the
fixed-point conditions counts every fixed configuration once.
Burnside averaging
\[
 \frac1{|\Gamma_K|}\sum_{g\in\Gamma_K}
 \bigl(\text{valid configuration series fixed by }g\bigr)
\]
therefore counts exactly the unlabelled graphs of this skeleton
and is rational.  Include the finite skeleton-label monomials
and the global factor $z^2$ from \eqref{eq:rational-node-weight},
sum over the finitely many skeletons, and add the complete-core
polynomial.  This proves Theorem~\ref{thm:fixed-excess-rationality}.

\paragraph{Effectivity and scope.}
All objects in this construction are finite and enumerable: bounded
labels and states, canonical incidence transitions, skeletons and
their automorphism groups, exact chain transformation maps,
regular-language intersections, and the half-word automata
\eqref{eq:rational-half-update}.  Weighted transfer matrices and
finite-group averages use exact rational arithmetic.  This gives
a terminating construction for any specified $e$, although its
transformation monoids can be very large.  Rationality implies
eventual constant-coefficient recurrences; it does not imply
polynomial growth, quasipolynomial coefficients, or only roots of
unity as poles.  Those stronger properties were proved explicitly
for $e=0,1$.  The present structural/automata proof has passed
independent mathematical audit; it is not an end-to-end Lean
formalization.
